\documentclass[10pt,a4paper]{amsart}
\usepackage[margin=19mm]{geometry}
\usepackage{amsmath,amssymb,mathtools}
\usepackage[scr=rsfs,cal=euler]{mathalfa}
\usepackage{xfrac}
\usepackage[unicode,hidelinks,bookmarks=false]{hyperref}
\newcommand{\ie}{i.e.\ }
\newcommand{\cf}{cf.\ }
\newcommand{\resp}{respectively\ }
\newcommand{\Subsection}{\subsection}
\providecommand{\nobreakdash}{\nobreak\hbox{-}}
\setlength{\emergencystretch}{2em}
\multlinegap0pt
\usepackage{amssymb}
\usepackage{xcolor}
\newtheorem{theorem}{Theorem}
\title{On Euler's magic matrices of sizes $3$ and $8$}
\author{Peter M\"uller}
\date{Source: arXiv:2504.16260v2 (CC BY 4.0)}
\begin{document}
\maketitle

\noindent\emph{Source and licence.} On Euler's magic matrices of sizes $3$ and $8$. Version arXiv:2504.16260v2 (CC BY 4.0), distributed by arXiv under the Creative Commons Attribution 4.0 International licence, http://creativecommons.org/licenses/by/4.0/, as recorded in the arXiv licence metadata for this exact version and shown on the abstract page https://arxiv.org/abs/2504.16260v2. Full licence notice: \texttt{licenses/arxiv-2504.16260-CC-BY-4.0.txt}.\par\smallskip

\noindent\textit{Editorial scope.} Counted unit: the article's theorem theorem (source0.tex lines 563--572). The article's complete original proof of it is delivered with it (source0.tex lines 573--586, one \texttt{proof} environment of 14 lines). No mathematics is written, repaired or re-derived: the statement and the proof are the article's own lines, in the article's own notation, and no retained line is substituted or re-flowed.  Retained uncounted as the source-local context the proof needs: lines 78--83.  Nothing was omitted from any retained span, so the package carries no omission record.  No retained line contains a citation, so the wrapper carries no bibliography and no bibliography entry.  No retained line contains a figure environment, an image-inclusion command or drawing code, so no image is delivered and the package declares no asset dependency.  Editorial formatting only, in addition: an AMS \texttt{amsart} preamble that restates the article's own declarations and macros used by the retained lines, namely the environments \texttt{theorem} on separate counters, together with the standard packages those lines need.  The retained environments print in this document as Theorem 1 in order of appearance, whereas the article prints the same items with the numbering of its own full document.  The complete native source of the article as distributed in the arXiv e-print for this exact version is bundled unchanged as \texttt{sources/176930/source0.tex}.
  \begin{align}
    \label{eq:1}
    M\cdot M^t &= \gamma\cdot I_n,\\
    \sum_{i=1}^nm_{i,i}^2 &= \gamma,\label{eq:2}\\
    \sum_{i=1}^nm_{i,n+1-i}^2 &= \gamma.\label{eq:3}
  \end{align}

\begin{theorem}
  Set \[ \sigma=
    \begin{cases}
      (1\;2\;\ldots\;n-1)(n),& \text{if }n\text{ is }even\\
      (1\;2\;\ldots\;k-1)(k)(k+1\;k+2\;\ldots n),& \text{if }n=2k-1%
                                                \text{ is odd.}\\
    \end{cases}
  \]
Then $M(\sigma)$ is an Euler's magic matrix.
\end{theorem}

\begin{proof}
  We need to verify the conditions \eqref{eq:1}, \eqref{eq:2}, and
  \eqref{eq:3}. As $M(\sigma)\cdot M(\sigma)^t=I_n$, condition
  \eqref{eq:1} holds with $\gamma=1$.
  
  First suppose that $n=2k\ge4$ is even. Then $m_{i,i}=0$ for
  $1\le i<n$ and $m_{n,n}=1$. Furthermore, $m_{i,n+1-i}=1$ if and only
  if $i=k$. We see that $M(\sigma)$ fulfills conditions \eqref{eq:2}
  and \eqref{eq:3}.

  Similarly, if $n=2k-1$ is odd, then $m_{k,k}=1$ and $m_{i,i}=0$ if
  $i\ne k$, and $m_{i, n+1-i}=1$ if and only if $i=k$. Again,
  $M(\sigma)$ fulfills conditions \eqref{eq:2} and \eqref{eq:3}.
\end{proof}

\end{document}
